\documentclass{article}
\usepackage{graphicx} % Required for inserting images

\title{Deep Learning: Final Project Report}
\author{Benjamin Cartwright }
\date{May 2026}

\begin{document}

\maketitle

\section{Introduction}

This project studies parameter-efficient adaptation of a vision-language model for science multiple-choice question answering from image-text context. Each example provides an image, a question, candidate options, and optional supporting context (hint and lecture). The goal is to predict the correct answer index.

Our base model is \texttt{HuggingFaceTB/SmolVLM-500M-Instruct}. We compare a zero-shot baseline against several adapter-based fine-tuning variants under a common prompt/evaluation setup. The main question is whether low-rank adaptation on selected transformer submodules can improve accuracy while keeping trainable parameters small.

The experimental strategy is:
\begin{itemize}
    \item establish a zero-shot reference point,
    \item train LoRA-based models with two objectives (option scoring and SFT),
    \item evaluate DoRA as a magnitude-direction variant of LoRA,
    \item test a broader LoRA target set that includes MLP and attention projections.
\end{itemize}

We report validation accuracy, compare adaptation choices, and analyze which parameter-efficient design provides the best performance/complexity trade-off.

\section{Dataset}

We use the provided CSV splits with paired images:
\begin{itemize}
    \item \texttt{train.csv}: 3109 examples,
    \item \texttt{val.csv}: 1048 examples,
    \item \texttt{test.csv}: 1008 examples.
\end{itemize}

Each row includes question text, a variable-length list of candidate choices, number of choices, image path, and (for train/val) the integer answer label. Additional context fields \texttt{hint} and \texttt{lecture} are included in most prompt variants. Images are loaded from split-specific directories under \texttt{data/images/images/\{split\}/...}.

Formally, each supervised sample is
\[
\left(x_i^{\text{text}},\, x_i^{\text{img}},\, \mathcal{C}_i=\{c_{i1},\dots,c_{iK_i}\},\, y_i\right),
\]
where $K_i$ is the number of options and $y_i \in \{1,\dots,K_i\}$ is the correct option index.

Minimal preprocessing is applied:
\begin{itemize}
    \item fill missing text fields with empty strings,
    \item parse serialized choice lists into ordered options,
    \item format prompts consistently (question + choices + optional hint/lecture),
    \item resize/process images through the model processor at runtime.
\end{itemize}

\section{Model Description}

Let $f_{\theta}$ denote the pretrained multimodal backbone, and let $p_{\theta}(a \mid x)$ be the model distribution over answer strings/tokens conditioned on image-text input $x$.

\subsection*{1) Zero-shot baseline}

No trainable parameters are added: $\theta$ is frozen. For each example we generate an answer string and parse it into an index:
\[
\hat{y} = \operatorname*{parse}\!\left(\arg\max_{a} p_{\theta}(a \mid x)\right).
\]
This provides a reference accuracy before adaptation.

\subsection*{2) LoRA model: option scoring and SFT}

For a linear projection $W \in \mathbb{R}^{d_{\text{out}}\times d_{\text{in}}}$, LoRA uses
\[
W' = W + \Delta W,\qquad
\Delta W = \frac{\alpha}{r}BA,
\]
with $A \in \mathbb{R}^{r\times d_{\text{in}}}$, $B \in \mathbb{R}^{d_{\text{out}}\times r}$, rank $r \ll \min(d_{\text{in}},d_{\text{out}})$, and frozen $W$.

\textbf{(a) Option scoring objective.}
For candidate answers $\{a_1,\dots,a_{K}\}$, define scores by conditional log-likelihood:
\[
s_k = \log p_{\theta,\phi}(a_k \mid x),
\]
where $\phi=\{A,B\}$ are trainable adapter parameters. Training minimizes cross-entropy over choices:
\[
\mathcal{L}_{\text{opt}} = -\log \frac{\exp(s_{y})}{\sum_{k=1}^{K}\exp(s_k)}.
\]

\textbf{(b) SFT objective.}
Given target answer token sequence $t_{1:T}$,
\[
\mathcal{L}_{\text{SFT}} = -\sum_{j=1}^{T}\log p_{\theta,\phi}(t_j \mid x, t_{<j}),
\]
with loss masking so only target-answer tokens are supervised.

\subsection*{3) DoRA model}

DoRA reparameterizes each adapted weight into magnitude and direction, updating direction with low rank while decoupling scale:
\[
W = m \,\hat{W},\qquad \|\hat{W}\|=1,\qquad
\hat{W}' = \hat{W} + \frac{\alpha}{r}BA,
\]
\[
W' = m' \,\frac{\hat{W}'}{\|\hat{W}'\|}.
\]
Compared with LoRA, this can improve stability by separating norm and directional updates.

\subsection*{4) LoRA with MLP + attention targets}

Standard LoRA often targets attention projections only:
\[
\mathcal{T}_{\text{attn}}=\{W_q,W_k,W_v,W_o\}.
\]
We also test broader coverage including feed-forward/MLP blocks:
\[
\mathcal{T}_{\text{attn+mlp}}=\{W_q,W_k,W_v,W_o,W_{\text{up}},W_{\text{down}},W_{\text{gate}}\}.
\]
Adapters are injected for each $W\in\mathcal{T}$:
\[
W' = W + \frac{\alpha}{r}B_WA_W.
\]
This increases adaptation capacity while remaining parameter-efficient; in experiments we compare whether this broader target set outperforms attention-only adaptation.

\end{document}
